\documentclass[11pt]{letter}
\usepackage[margin=1in]{geometry}
\usepackage{parskip}
\usepackage{hyperref}

\signature{Ning Zhou\\
Public R\&D Center, SuperMap Software Co., Ltd.\\
Beijing, China\\
\href{mailto:zhouning1@supermap.com}{zhouning1@supermap.com}}

\address{}

\begin{document}

\begin{letter}{The Editor \\ \textit{Nature Computational Science}}

\opening{Dear Editor,}

We are pleased to submit our manuscript, ``Pairwise ranking restores model-based planning in high-branching discrete action spaces, with application to county-scale farmland consolidation,'' for consideration as a Research Article in \textit{Nature Computational Science}.

Model-based reinforcement learning is widely viewed as the natural fit for combinatorial planning problems with structured state and a known cost function, yet applications to high-branching discrete action spaces have repeatedly under-performed model-free baselines without a clear diagnosis. Our manuscript identifies why and shows the fix. We believe four aspects of the work align with the journal's interests in methodological clarity, computational rigour, and reproducibility.

First, the methodological diagnostic. We show that aggregate prediction fidelity (cosine similarity $0.998$) and within-state action ranking accuracy ($51.6\%$, statistically indistinguishable from random) are independent quantities under mean-squared-error training, and that the latter is what model-predictive top-$K$ planning actually consumes. The decoupling is established empirically across a $\sim$$2{,}600$-way discrete action space, supported by an analytic argument (the conditional-mean predictor minimises MSE while giving $50\%$ rank accuracy), and shown to account for four otherwise-puzzling failure modes of the learned planner. The implication for evaluation methodology is that learned models passing standard aggregate-fidelity tests on continuous-control or low-branching benchmarks may fail catastrophically when transferred to high-branching discrete domains, a regime currently under-represented in the model-based RL literature.

Second, the fix. Augmenting the existing reward head with a pairwise large-margin loss (a classical learning-to-rank technique) raises ranking accuracy to $85.5\%$ without architectural change, and converts model-predictive control from the worst variant tested to the strongest. The intervention is conceptually adjacent to MuZero's policy/value scorers, AlphaZero-style move evaluators, and RLHF reward modelling---all discriminative learned-scoring approaches---but applied at the model side of the planning pipeline rather than the policy side, in a regime where the generative latent-dynamics framing of contemporary world models brings no benefit and the planner only needs the model to rank, not to roll out.

Third, the verification. We address head-on the question of whether reported gains are real geometry or learned-dynamics artefacts. The manuscript reports a five-layer verification stack including: (i)~independent GIS recomputation of all three primary metrics directly from the optimised cadastral shapefile, agreeing with planner-reported numbers to floating-point noise without invoking any learned component; (ii)~ensemble-member subset ablation establishing operational variance under the deterministic pipeline; (iii)~direct measurement of the ensemble's 1-step prediction error against the true environment along the same trajectory MPC visits (reward MAE $0.82$ against true-reward standard deviation $1.40$, yet Spearman rank correlation $+0.22$ sufficient at $H=5$); (iv)~a counter-factual ablation replacing the ensemble with the perfect dynamics model that is strictly worse on all three objectives; and (v)~cross-county replication on a second region of the same regulatory regime. To our knowledge this is the deepest verification stack offered in the published RL-for-spatial-planning literature.

Fourth, the reproducibility infrastructure. We release the trained ensembles ($48$ checkpoints across two counties and a $\lambda$-ablation set), a deterministic seven-preset synthetic benchmark under CC-BY 4.0, an end-to-end reproduction recipe walking from raw DLTB shapefiles plus public Copernicus DEM tiles to the §5 cross-seed numbers in $\sim$$2$ hours per county, and per-county reproduction artefacts whose schema diffs field-by-field against the research-side originals. An external reader can verify the pairwise-ranking diagnostic on user-supplied data using the released evaluation script, and an external county planner can run the deployment toolbox without access to the original cadastral data via the bundled synthetic landscapes.

We further release a deployable ArcGIS Pro Python Toolbox bringing the full pipeline within reach of GIS practitioners without machine-learning support, with a 70-second continuous-integration smoke test and approximately seven minutes per 100-step planning scenario on a desktop CPU.

The manuscript has not been submitted, in whole or in part, to any other journal. The authors have no competing interests. All data presented are original; no identifying cadastral records are redistributed; the released benchmark contains only synthetic landscapes from a deterministic seed-controlled generator. We thank the Editor and reviewers for their time and look forward to their assessment.

\closing{Sincerely,}

\end{letter}
\end{document}
